% !TEX root = ../main.tex Preamble %%% ------------------------------------------------

% ## 1.0: Preamble

% P1: BSP Is Important

% P2: Data + Compute challenges

\chapter{Introduction}
\minitoc

The Bulk Synchronous Parallel (BSP) paradigm powers some of our most critical computational tools,
from large-scale scientific simulations in astrophysics and climate science to distributed training
of foundation models in machine learning.
These workloads---consisting of large parallel computations interspersed with I/O and global
synchronization---currently operate at the scale of tens of thousands of accelerator
nodes~\cite{Das2023:LargeScale,Stock2024:LargeScale,Gratt2024:Llama3Scale,:LlamaScale} and
continue to grow with time~\cite{Rahma2024:MLScale}.
As these applications scale, their efficiency has datacenter-level impacts on both cost and
time-to-discovery.

The tightly synchronized nature of BSP creates distinct challenges along the data and compute
paths.
On the data path, bursts of checkpointed state create massive I/O demands on the storage and
analytics
infrastructure~\cite{Jones2012:Checkpointing,Eisen2022:Checkpoint,Lian2025:UniversalCheckpointing}.
On the compute path, frequent global synchronization amplifies the system's sensitivity to
execution
variability~\cite{Hoefl2010:Noise,Petri2003:Missing,Yildi2016:IONoise,Kocol2018:Varbench,Acun2016:Variation},
and also amplifies correctness risks when failures or corrupted state propagate across ranks.
Runtime artifacts that would be manageable in loosely coupled systems, such as minor OS-level
interruptions or thermal throttling, become critical scaling barriers when every worker must wait
for the slowest participant.

Across two studies spanning the data and compute paths, we find that many of these efficiency
bottlenecks are driven by \emph{runtime unknowns}: conditions that cannot be predicted reliably
before execution.
In this regime, the decisive factor is how quickly these conditions can be observed, characterized,
and mitigated.
On the data path, CARP tackles range-query performance under skewed and evolving key distributions
by reconstructing global distributions in situ and adapting partitioning as the run evolves.
On the compute path, our placement study in \emph{adaptive mesh refinement} (AMR) codes shows
the failure modes of slow turnaround:
placement effects were repeatedly confounded by fail-slow hardware, fabric
pathologies, and tuning artifacts, forcing months of reruns and offline telemetry analysis.
We view these as two instances of the same primitive---a runtime feedback loop over distributed
state---online and closed in one case, manual and slow in the other.
Because the governing conditions are fundamentally unknowable in advance, addressing them in
production requires general-purpose infrastructure for closing these loops,
which BSP workloads largely lack today.

BSP's tight coupling is both the source of these bottlenecks and what makes them structurally
addressable.
Closing a runtime feedback loop requires two capabilities: observation to materialize a view of
distributed execution, and action to respond.
The same synchronization boundaries that amplify every rank's slowdown into a system-wide stall
also define natural points where both are meaningful.
We call the resulting paradigm steerable observability: the ability to navigate across views at
multiple fidelities during the run and to intervene on the application guided by those views.
We propose ORCA (\emph{Observability with Real-time Control and Aggregation}), a programmable
substrate for feedback and control aligned to BSP synchronization,
so the loop operates on well-defined boundary snapshots rather than devolving into reruns and
post-hoc pipelines.

\begin{tcolorbox}[thesisbox]
  \textbf{Thesis Statement:} A large class of efficiency bottlenecks in BSP systems arise from
  conditions that are unknowable before runtime.
  Through two studies spanning the data and compute paths, this thesis demonstrates that such
  conditions are systematically addressable through adaptive mechanisms if the requisite feedback
  loops exist, and perpetuate if they do not.
  We present ORCA as general-purpose infrastructure for closing these loops at scale.
\end{tcolorbox}

The rest of this chapter makes the case for programmable feedback loops through two studies
spanning the data and compute paths (\S\ref{sec:intro:case}), presents steerable observability and
its
realization in ORCA (\S\ref{sec:intro:orca}), lists contributions (\S\ref{sec:intro:contribs}), and
outlines the rest of the
thesis (\S\ref{sec:intro:org}).

\section{The Case for Programmable Feedback Loops in BSP}
\label{sec:intro:case}

We begin by defining the anatomy of a feedback loop over BSP execution
(\S\ref{sec:intro:case:anatomy}), then examine
two studies that instantiate the pattern with different outcomes.
CARP demonstrates a fixed-function
online loop for adaptive range partitioning on the data path (\S\ref{sec:intro:case:carp}).
Our AMR placement study
demonstrates the cost when the loop is manual and offline on the compute path
(\S\ref{sec:intro:case:amr}).
We
conclude by characterizing the infrastructure gap that separates the two
(\S\ref{sec:intro:case:gap}).

\subsection{Anatomy of a BSP Feedback Loop}
\label{sec:intro:case:anatomy}

BSP execution is computation punctuated by synchronization.
Ranks compute independently, then coordinate at barriers and collectives.
These boundaries define when distributed execution is globally interpretable and when coordinated
change is well-defined.
A feedback loop over BSP execution operates on this cadence.

\paragraph{View} A view is a reduction over per-rank observations at a synchronization boundary
that answers a specific question---what is the load distribution across ranks, where is the
synchronization bottleneck, what is the I/O throughput at this checkpoint.
The point is to materialize the view that makes the condition legible, not to export raw telemetry.

\paragraph{Action}
An action is a coordinated update applied at a synchronization boundary so all ranks remain
consistent.
One class of actions steers observability itself, changing what is observed, at what fidelity, and
where, so subsequent views answer the next question.
A second class steers the running application or runtime, applying boundary-consistent updates such
as hyperparameter changes, repartitioning, or placement adjustments.

\paragraph{Operator}
The operator requests views and chooses actions.
It may be a fixed policy, a human workflow, or an automated agent.
It is external to the substrate.
The substrate provides boundary-consistent observation and action.
The operator supplies the policy that closes the loop.


\subsection{CARP: Adaptive Online Range Partitioning}
\label{sec:intro:case:carp}

Range queries over scientific data are fundamental to analysis, but expensive to support at scale.
Post-processing approaches like out-of-core sorts require multiple additional passes over data,
incurring write amplification and reducing the effective application bandwidth.
While in-situ hash partitioning for point lookups has been demonstrated viable for scientific data,
range partitioning is harder --- key distributions are unknown, highly skewed, and evolve over
time.

CARP addresses this with adaptive online range partitioning driven by periodic reconstruction of
the global key distribution.
A \emph{renegotiation} protocol materializes the global key
distribution via in-situ reductions over per-rank state when triggered.
A policy triggers this protocol at a preconfigured intervals to obtain the current
global distribution, which is then used to compute new load-balanced partition boundaries and drive
shuffling en-route to storage.
By range-partitioning data in a single streaming pass,
CARP avoids the write amplification of post-hoc sorting while
enabling comparable query performance, achieving up to $4.9\times$ faster ingestion in our
evaluation.

\subsection{AMR Placement: A Manual Loop for Compute Pathologies}
\label{sec:intro:case:amr}

\emph{Adaptive Mesh Refinement} (AMR)
codes are widely considered challenging to load-balance.
The mesh evolves dynamically to track changing physics, requiring frequent rebalancing, but
load characteristics and placement effects have not been well studied in practice.
We initially set out to build telemetry-driven placement policies that could adapt to these
characteristics at runtime.
In practice, we found that we could not even evaluate placement cleanly: placement effects were
repeatedly confounded by unrelated artifacts—fail-slow hardware, fabric pathologies, and tuning
issues—that had to be identified and eliminated before placement could be isolated.

The resulting loop---collect traces, analyze offline, refine hypothesis, reconfigure and
rerun---was entirely manual and slow.
Conventional HPC tracers emit large per-rank logs in formats
that required expensive \emph{extract-transform-load} (ETL) before any analysis could begin.
% New hypotheses would require repeating the loop with a reconfigured setup.
Eliminating confounders consumed months of iteration even when individual mitigations typically
took days once identified.
With a stable measurement baseline finally established, we found a
fundamental tradeoff between compute load balance and communication locality, and designed a
placement policy called CPLX to provide tunable control over this tradeoff.
In evaluation, CPLX reduced runtime by up to 21.6\% in real AMR workloads over tuned baselines, but
like
CARP's renegotiation interval, exposed another context-dependent tunable knob.

\subsection{The Infrastructure Gap}
\label{sec:intro:case:gap}

Both studies instantiate the same primitive---a feedback loop over distributed execution---but
with starkly different costs.
In CARP, the loop is online: it materializes the specific view it
needs during the run and uses it to drive reconfiguration with low turnaround.
The AMR loop was expensive, as telemetry was collected as bulk traces and analyzed offline.
Offline loops are expensive in three compounding ways:

\begin{itemize}[leftmargin=*]
  \item \textbf{Deferred Materialization Costs.
        }
        Raw telemetry is captured at full volume and written
        to storage before any reduction occurs.
        Because the right subset is not knowable in advance,
        the paradigm intrinsically tends to overcollect.
        The resulting dataset is supercomputer-scale,
        so a reduction that could run in situ over the fabric is instead forced through a
        write--read--parse cycle on the slowest tier in the hierarchy: storage.

        % \item \textbf{Deferred Materialization Costs.} Raw telemetry is captured at full volume,
        % serialized, and written to storage before any reduction occurs. Because the right subset is not
        % knowable in advance, the paradigm intrinsically tends to overcollect. The resulting datasets are
        % cluster-scale, and analyzing them demands the same computational resources that in-situ processing
        % would leverage---but mediated by the slowest tier in the hierarchy (storage).

  \item \textbf{ETL Costs.
        }
        Conventional tracers write per-rank logs in flat layouts
        with little awareness of BSP structure.
        Causal relationships created by synchronization are not reflected in the on-disk organization,
        so even simple questions require expensive ETL before queries can run.

  \item \textbf{Rerun Costs.
        }
        Refining a hypothesis, changing what is observed, or
        testing a mitigation typically requires reconfiguration and re-execution.
        Turnaround is bounded by
        reruns and offline pipelines rather than by the running job.
\end{itemize}



The real cost is foregone analysis: when turnaround is dominated by trace volume, ETL, and reruns,
systematic investigation gives way to ad-hoc mitigations.
% The real cost is not just wasted cycles, but foregone analysis. When turnaround is dominated by
% trace volume, ETL, and reruns, many questions are never investigated systematically, and teams
% fall back to ad-hoc mitigations that treat symptoms rather than causes. CARP shows that online
% closure is feasible, but only as a fixed-function mechanism built for one view and one action, not
% as general-purpose infrastructure that lets operators specify new views and interventions. These
% are not pathologies unique to our studies. They recur across BSP systems and are extensively
% documented, yet persist because BSP workloads largely lack general-purpose infrastructure for
% closing such loops online.
CARP shows that online closure is feasible, but only as a fixed-function mechanism built for one
view and one action, not as general-purpose infrastructure that lets operators specify new views
and interventions.
These pathologies recur across BSP systems and are well-documented, yet persist because no
general-purpose substrate for closed-loop observability workflows exists.

\section{ORCA: A Substrate for Closed-Loop Observability}
\label{sec:intro:orca}

The preceding analysis points to a missing piece: a programmable substrate for closing feedback
loops over BSP execution online, supporting operator-specified views and coordinated interventions
without bespoke pipelines.
We call the resulting paradigm steerable observability.
The system materializes views of distributed state on demand, operators navigate across fidelities
at runtime, and interventions are applied consistently at synchronization boundaries guided by
those views.
This makes always-on collection viable: a lightweight summary can run continuously, with detail
enabled only when needed.

\paragraph{Requirements} We identify three requirements for steerable observability:

\begin{itemize}
  \item \textbf{In-situ view materialization.
        }
        Views must be reduced in-situ during the run, not
        deferred to post-hoc reconstruction through the storage tier.
        Overhead should scale with fidelity:
        always-on lightweight aggregates in the common case, with detailed collection enabled on demand.

  \item \textbf{Programmable views.
        }
        Operators must be able to specify and evolve views as queries
        over per-rank observations, without recompilation, restarts, or bespoke pipelines---not
        write-optimized per-rank logs in flat formats that require expensive reshaping before any question
        can be asked.
        A fixed catalog of views cannot anticipate every pathology.

  \item \textbf{Control semantics.
        }
        A well-defined mechanism for applying coordinated interventions
        at synchronization boundaries, so that refining a hypothesis or changing what is observed does not
        require reconfiguration and re-execution.
\end{itemize}

\paragraph{ORCA} We propose a \emph{tree-based overlay network} (TBON) called ORCA
(\emph{Observability with
  Real-time Control and Aggregation}) as a programmable substrate for steerable observability in BSP
workloads.
ORCA attaches a tier of dedicated aggregators to a running application and executes SQL
queries over columnar telemetry in the overlay, materializing boundary-aligned views through
in-situ reduction.
A timestep-consistent control plane (TS2PC) applies coordinated interventions,
including changes to what is observed, atomically at barriers.
On a real AMR code at 4096 ranks,
ORCA collects detailed traces at 1--2\% overhead versus 56--81\% for conventional tracers,
eliminates 98.5--99.999\% of telemetry at source for evaluated queries, and in a closed-loop
evaluation with an LLM agent localizes a performance anomaly in 27 minutes that previously took
months of manual iteration.

\paragraph{Scope}
ORCA provides two capabilities: in-situ view materialization and timestep-consistent control.
ORCA is agnostic to how telemetry is produced and to who closes the loop---it may be steered by a
fixed heuristic, a human workflow, or an automated agent.
Telemetry can originate from probes, counters, or application annotations, as long as it can be
represented as structured columnar data.
Where SQL is insufficient, ORCA supports user-defined operators within the overlay.


\section{List of Contributions}
\label{sec:intro:contribs}

\paragraph{CARP}
\begin{itemize}
  \item An adaptive in-situ range partitioner for streaming scientific data, designed to accelerate
        range queries under unknown and evolving key distributions.

  \item Partitions data in a single streaming pass with $1\times$ write amplification during
        ingestion,
        and is $2.8$--$4.9\times$ faster than post-processing approaches.

  \item Shows that partially ordered layouts can deliver range-query performance comparable to fully
        sorted layouts, while being up to two orders of magnitude faster than auxiliary bitmap indices
        on the evaluated workloads.

  \item Complements DeltaFS's in-situ hash partitioning~\cite{Zheng2018:DeltaFS} by providing the
        in-situ range-partitioning
        counterpart needed for range queries.
\end{itemize}

\paragraph{AMR placement}
\begin{itemize}
  \item An end-to-end empirical study integrating collection, tuning, and placement, showing how
        artifacts such as overheating, fabric bugs, and operation ordering obscure placement effects.

  \item CPLX, a tunable placement policy that exposes the tradeoff between compute load balance and
        communication locality, achieving up to $21.6\%$ runtime reduction over tuned baselines across
        $512$--$4096$ ranks.
\end{itemize}

\paragraph{ORCA}
\begin{itemize}
  \item A BSP-aware observability and control substrate that aligns both analysis and coordinated
        interventions to synchronization boundaries.

  \item The \emph{timestep dataframe} abstraction, partitioning BSP telemetry into causally
        independent windows,
        and a \emph{(timestep, rank)}-partitioned
        Parquet layout, yielding up to $44\times$
        smaller traces and up to $6{,}800\times$ faster post-mortem analysis.
        % grounded in the fact that global synchronization% partitions execution into causally independent
        % windows, 

  \item An incast-optimized RDMA transport for tree topologies that enables detailed trace collection
        at under $2\%$ runtime overhead with real AMR codes, versus $56$--$81\%$ for conventional tracers.

  \item A SQL-based in-situ programming interface over timestep dataframes with operator pushdown,
        eliminating $98.5\%$--$99.999\%$ of telemetry
        at source in evaluated queries.

  \item A timestep-consistent control plane via TS2PC (\emph{timestep-linked two-phase commit}) that
        enables coordinated runtime changes aligned to BSP synchronization.

  \item In a closed-loop evaluation using an LLM agent, an anomaly that took months to isolate
        manually is localized in 27 minutes.
\end{itemize}


\section{Thesis Organization}
\label{sec:intro:org}

The remainder of this thesis is organized as follows.
\begin{itemize}
  \item \cref{chap:bg} provides background on BSP environments and efficiency bottlenecks.
  \item \cref{chap:carp} presents CARP and its evaluation on ingestion/query performance.
  \item \cref{chap:amr} presents the AMR diagnostic workflow and CPLX placement policy.
  \item \cref{chap:orca} presents ORCA and its observability/control design.
  \item \cref{chap:timeline} outlines remaining milestones and completion plan.
\end{itemize}

% P3: We studied this and found a common shape: "runtime unknowns" + "adaptive feedback loops".
% Efficiency challenges of the "runtime unknown" form become hard to address in prod because the
% general shape in a necessary form is not supported.

% P4: BSP is unique, this shape needs to be solved in a unique way. ORCA does it.

% P5: Thesis Statement

% P6: Intro Outline

% ## 1.1: Feedback Loops + Gaps

% Takeaway: When addressing runtime unknowns, BSP loops keep showing up. Yet no general mechanism
% exists to build these loops.

% ### 1.1.1: Loops and BSP

% BSP execution is work punctuated by sync. This sync is both the primary causality and consistency
% mechanism

% State: System State. Global state only makes sense at sync points. Local state/internal mutations
% are a second-order explanation. - This holds true for both correctness and performance

% View: Aggregates from that state. The right aggregate makes problems obvious. Reductions. (Not
% necessarily in-situ)

% Action: application-wide actions only defined at these synchronization points.

% Operator: holds the policy. Decoupled from the substrate. Can be algorithm or human.

% ### 1.1.2: CARP: Functional, Fixed-Function, Data Path Loop, Driven By Policy

% ### 1.1.3: AMR: Broken, Generic, Driven By Human

% ### 1.1.4: The Gap

% 1. Delaying materialization makes it costly and derails interventions - You overcollect,
% serialize, exit a HPC fabric, write to storage. - Because this is the state from a supercomputer,
% reversing this needs supercomputer-sized resources. And you have added the slowest tier on the
% path (storage). - You overcollect because you don't know how much to collect - Or because you're
% collecting raw data. Out-of-core sorts first collect raw data, look at histograms, then move it
% around. Couldn't you look at histograms and organize it properly in the first place?

% 2. Given a certain collection volume, we do a terrible job organizing it - Write JSON, with no
% understanding of these causal boundaries. Per-rank logs are write-optimized they're not
% structure-optimized. - There's an RDMA fabric but we don't use it

% 3. Interventions require trial and error. "What's the right value of this hyperparameter"? "How
% many probes to toggle? There's infinite." Each trial and error becomes a full run.

% %%% - In-situ reductions. Views are simple aggregates. Collecting data needlessly adds overhead.
% If you can reduce in-situ you can act in-situ.


% ## 1.2: Steerable Observability and ORCA

% Paradigm: Steerable and view-driven observability. System materializes what you specify on demand,
% and enables you to act in real time. This forms a complete feedback and control loop. Makes
% always-on possible as there's always a meaningful summary you're acting on. Details on demand.

% ### Steerability Needs: - Fast views materialized in-situ. - Programmability: views vary - Online
% intervention: reduce the cost of a trial

% ### How ORCA Saves The Day - TBON and all

% ### What Is And Not In Scope - In-scope: a SQL-based substrate + timestep-consistent control - Not
% in scope: instrumentation, or policies - These are context-dependent. Decoupling is the right
% thing to do - Our goal: make it easy to add these things. Anything that materializes Arrow is in
% scope - Likewise for policies. Prior work does autotuning. For us that is out of scope. -
% Operators. Many things SQL doesn't enable easily. Specific kernels may be needed. We acknowledge
% that and provide an escape hatch.

% 1.3: List of Contributions

% 1.4: Thesis Outline
